\documentclass[11pt,a4paper]{article}
\usepackage[utf8]{inputenc}
\usepackage[T1]{fontenc}
\usepackage{lmodern}
\usepackage{textcomp}
\usepackage{amsmath,amssymb}
\usepackage{graphicx}
\usepackage{booktabs}
\usepackage{longtable}
\usepackage{array}
\usepackage{listings}
\usepackage{xcolor}
\usepackage[margin=1in]{geometry}
\usepackage{float}
\usepackage[hidelinks,breaklinks]{hyperref}
\usepackage{microtype}
\setlength{\parskip}{4pt}
\lstset{basicstyle=\ttfamily\scriptsize,breaklines=true,breakatwhitespace=false,
  postbreak=\mbox{\textcolor{gray}{$\hookrightarrow$}\space},frame=single,
  framerule=0.3pt,rulecolor=\color{gray!60},columns=fullflexible,keepspaces=true,
  showstringspaces=false,numbers=none,xleftmargin=4pt,xrightmargin=4pt,
  literate={\_}{{\_}}1}
\graphicspath{{figures/}}
\providecommand{\tightlist}{}
\title{The 1.0--1.2 TeV Thermal Higgsino Against Planck, LZ and XENONnT: Toy-Model Relic and Inelastic-Recoil Calculations with a Literature Survey}
\author{HEP-Genesis-Agent session record \\ \small Project: ramp-combo-higgsino-dd}
\date{2026-10-22}
\begin{document}
\maketitle
\noindent\textit{\small This document is an AI-generated account of an interactive session with the HEP-Genesis agent. Its text was written by a language model strictly from the session record; figures and numbers come only from tool results obtained during the session. The supplementary material and the transcript, generated alongside, document exactly what ran. Verify independently before use.}
\vspace{1em}

\begin{abstract}
We report a two-pass investigation of where the pure higgsino thermal relic, at $m_\chi \approx 1.0$--$1.2$ TeV, sits relative to the Planck abundance and to current xenon direct-detection limits. We combined calculations from an inelastic-dark-matter (inelastic-DM) tool server with an abstract- and HTML-level literature survey.

The server's relic model is a single-channel Kolb--Turner freeze-out with a built-in higgsino cross section $\langle\sigma v\rangle = 2.5\times10^{-26}\,(1100\,\mathrm{GeV}/m)^2\ \mathrm{cm^3\,s^{-1}}$. It gives $\Omega h^2 = 0.0731$, $0.0898$, $0.1083$, $0.1285$ and $0.1503$ at $m = 0.9$--$1.3$ TeV, and crosses the Planck value $\Omega_c h^2 = 0.120\pm0.001$ at $m \simeq 1159$--$1160$ GeV. That crossing reflects the calibration and is not a Sommerfeld or coannihilation prediction. The cited literature places the thermal mass at about 1.08--1.2 TeV.

In the tool, elastic $Z$-mediated scattering at zero splitting predicts $\sim1.1\times10^{7}$ LZ events and is excluded. With the tool's hard halo cutoff, $v_{\rm esc}+v_E = 776$ km/s, about one event in the LZ 2.84 t$\cdot$yr extended window occurs only near $\delta \approx 352$ keV, and the rate closes by 370 keV. The literature quotes an IceCube solar-capture bound $\delta > 566$ keV, which lies beyond xenon's inelastic reach.

For elastic spin-independent scattering, the pure higgsino is only bounded, at $\sigma_{\rm SI}\lesssim10^{-48}\,\mathrm{cm^2}$. At 1.1 TeV this is at least $\sim33\times$ below the LZ 4.2 t$\cdot$yr limit, at least $\sim400\times$ below the XENONnT 3.1 t$\cdot$yr limit, and at least $\sim5.5\times$ below the xenon neutrino floor.

Several tool flags proved unreliable, halo dependence could not be tested, and indirect-detection conclusions rest on abstracts only.
\end{abstract}

\section{Introduction}
The user asked four things. First, a survey of recent arXiv work on the higgsino thermal relic target at 1.0--1.2 TeV. Second, use of the available inelastic-DM tools to compute the relic density at several masses. Third, a comparison with LZ/XENONnT sensitivities and the Planck abundance. Fourth, a statement of where the higgsino sits relative to current limits. The same request was issued twice. The first pass ran with a largely non-functional arXiv API layer. The second pass, after the API recovered, added verbatim citations, elastic spin-independent (SI) limits at TeV masses, a relic cross-section sensitivity test, and a gaugino-scale translation of the mass splitting.

The research plan evolved through the following sub-questions:
\begin{enumerate}
\item The tool's relic curve and its Planck crossing.
\item The inelastic LZ event counts and the kinematic edge of xenon.
\item Verbatim literature support for the key numbers.
\item The width of the splitting window that gives about one LZ event, and its halo sensitivity.
\item The indirect-detection status of the higgsino, added because the server's annihilation treatment was judged physically wrong.
\item Elastic SI placement relative to limits and the neutrino fog.
\item The gaugino-mass scale implied by the splitting.
\end{enumerate}

Much of the 2026-dated literature concerns a single LZ nuclear-recoil event at 248 keV (arXiv:2609.02823). It enters here only insofar as it bears on the higgsino's position relative to limits.

\section{Methods}
\paragraph{Tool server.} Calculations used an \texttt{inelastic-dm-server} with local dispatch. Its self-description cites the Lewin--Smith rate formula, Kolb--Turner freeze-out, Fermi-LAT limits, and Planck 2018 for CMB energy injection. It stores $\Omega_{\rm DM}h^2 = 0.12\pm0.0012$, the LZ event at $248\pm23$ keV, and an exposure of 2.84 t$\cdot$yr. The tools used were:
\begin{itemize}
\item \texttt{compute\_relic\_density}, with inputs $m_\chi$, an \texttt{is\_higgsino} flag and an optional $\langle\sigma v\rangle$;
\item \texttt{predict\_recoil}, with inputs $m_\chi$, $\delta$, $\sigma_n$ and the higgsino flag. It returns expected counts in the LZ extended nuclear-recoil window, $v_{\min}$ at the spectrum peak and at 248 keV, and an accessibility flag;
\item \texttt{check\_all\_constraints}, covering relic, LZ, Fermi-LAT and CMB $p_{\rm ann}$;
\item \texttt{scan\_mass\_splitting} and \texttt{plot\_constraint\_summary};
\item \texttt{predict\_annual\_modulation} and \texttt{plot\_recoil\_spectrum}.
\end{itemize}
With the higgsino flag set, the tool assigns $\sigma_n = 7.43\times10^{-39}\,\mathrm{cm^2}$ per nucleon, corresponding to $Z$ exchange. None of the tools accepts halo parameters. The halo is hard-coded with $v_{\rm esc}+v_E = 776$ km/s.

\paragraph{Runs.} The relic density was computed at $m = 900$, 1000, 1100, 1200, 1300 and 1159 GeV. Recoil predictions at $m = 1100$ GeV covered $\delta = 0$, 300, 330, 350, 355, 360, 365, 370, 375, 380, 390, 500, 566 and 700 keV. Further recoil runs were made at $m = 1000$ and 1200 GeV ($\delta = 380$ keV), at 1159 GeV ($\delta = 360$ keV), and at 1160 GeV ($\delta = 350$ and 360 keV).

Constraint checks were run at:
\begin{itemize}
\item $m = 1100$ GeV with $\delta = 0$, 200, 350, 400, 566 and 1000 keV;
\item $m = 1159$ GeV with $\delta = 400$ keV;
\item $m = 1080$, 1160 and 1200 GeV with $\delta = 350$ keV.
\end{itemize}

Two parameter scans were made: $17\times17$ over 800--1400 GeV $\times$ 0--800 keV, and $15\times15$ over 1000--1200 GeV $\times$ 0--700 keV. The second wrote to the same output file as the first. To probe the relic normalization, $\langle\sigma v\rangle$ was varied at 1100 GeV ($2.0$, $2.25$, $3.0\times10^{-26}$ cm$^3$/s) with the higgsino flag off. This was necessary because with the flag on the explicit argument was ignored and the built-in value was echoed back.

\paragraph{Hand calculations.} The kinematic edge was estimated by hand from $\delta_{\max}=\mu v_{\max}^2/2$, with $m_{\rm Xe}\approx122$ GeV and $\mu=109.8$ GeV at 1.1 TeV. The gaugino scale was estimated from the tree-level relation quoted from arXiv:1410.4549.

\paragraph{Literature access.} In the first pass, the direct arXiv API search returned HTTP 503/429 errors. Metadata, download and \texttt{cite} calls returned ``not found'' even for arXiv:1807.06209, and the full-paper downloader failed with \texttt{'Search' object has no attribute 'results'}. Only an OpenAlex-backed search returning abstracts worked.

In the second pass, metadata lookups succeeded and \texttt{cite} recorded verbatim abstract quotes for five papers. A sub-agent read arXiv HTML pages (downloads still failed), the LZ HEPData record 155182 (table t1), and O'Hare's public neutrino-floor file \texttt{Xe\_SI.txt}. A \texttt{cite} attempt on arXiv:2410.17036 failed its quote check.

Figure extraction from arXiv:2410.17036 and arXiv:2502.18005 returned only event-distribution and calibration plots, not the SI exclusion curves. These figures carry no information used in the analysis and are not shown. The full tool-call log and code are in the supplement.

\paragraph{Reporting discrepancy.} In both passes, the assistant's summary text claimed that no inelastic-DM or relic-density tools were available and that no relic densities were computed. The tool record contradicts this: the server tools were called and returned results. The results below are taken from those tool outputs.

\section{Results}
\subsection{Literature survey of the thermal target}
All values in this subsection come from abstracts. The 2026 identifiers were first found only through OpenAlex and were later confirmed by arXiv metadata in the second pass, at least for arXiv:2609.02775 and arXiv:2609.02608.

The verbatim-cited mass values are as follows:
\begin{itemize}
\item arXiv:2410.15733 states that in the thermal-relic scenario ``the observed DM abundance is realized with the Higgsino mass of $\sim$1.2 TeV''. The abstract says it includes coannihilation and the Sommerfeld effect. The paper itself is about non-thermal production; the thermal value is quoted there for contrast.
\item arXiv:2609.02775 uses ``the cosmologically preferred mass, $m_\chi\simeq1.08$ TeV''.
\item arXiv:2609.02608 states that the relic abundance ``fixes the mass near 1.1 TeV''.
\item arXiv:2403.19598 states that a thermal higgsino-like state ``must be between about 100 and 1150 GeV''.
\end{itemize}
The Planck 2018 abstract (arXiv:1807.06209) gives $\Omega_c h^2 = 0.120\pm0.001$.

Only arXiv:2410.15733 states its Sommerfeld and coannihilation treatment, so the cause of the 1.08 versus 1.2 TeV spread was not established. A ``1.08--1.2 TeV band with central value near 1.1 TeV'' is an inference from the spread of quoted values, not a published result. The assistant also relayed, from abstracts, that lighter thermal higgsinos are under-abundant by a factor of 5--20 in natural SUSY (arXiv:1803.11210).

For the LZ event, the extended-window search (arXiv:2609.02823) reports an exposure of 2.84 t$\cdot$yr, a window up to about 270 keV, and one event at $248\pm23\,(\mathrm{stat})\pm23\,(\mathrm{sys})$ keV. Its global significance is $2.6\sigma$ and its maximum local significance $3.4\sigma$. In the first pass these quotes came from OpenAlex and were marked contested.

Preferred splittings quoted in the literature are:
\begin{itemize}
\item about 350 keV for the Standard Halo Model and about 500 keV with Large Magellanic Cloud effects (arXiv:2609.01504);
\item $\delta\simeq377$ keV (arXiv:2609.02608);
\item 481--519 keV for $m_\chi=1000$--1183 GeV including the LMC tail (arXiv:2609.01590). The same abstract reports that ten years of IceCube data exclude $\delta\lesssim475$ keV near elastic blind spots;
\item a sideband optimum of 339 keV at 1 TeV (arXiv:2609.26698, as relayed by the assistant).
\end{itemize}

The decisive constraint is the IceCube solar-capture bound, $\delta>566$ keV at $m_\chi\simeq1.08$ TeV (arXiv:2609.02775), verbatim cited in the second pass. That abstract states that it ``exclude[s] the interpretation of the recent LZ event in terms of endothermic inelastic scattering of Higgsino dark matter''. The assistant additionally relayed, without verbatim quotes, the following:
\begin{itemize}
\item supporting solar-capture results (arXiv:2609.11833, arXiv:2609.07807, and 0.51--0.56 MeV in arXiv:2609.06760);
\item sideband tension (arXiv:2609.09385, arXiv:2609.26698);
\item the claim that halo uncertainties cannot close the gap (arXiv:2609.15321);
\item a surviving Planck-compatible branch at $\delta\approx0.85$--1.34 MeV (arXiv:2610.00226).
\end{itemize}
The abstract of arXiv:2609.09136 states that the thermal higgsino with $\delta\sim350$ keV is ``essentially excluded'' and proposes a non-higgsino inelastic candidate.

\subsection{Relic density from the tool}
Table~\ref{tab:relic} lists the tool's relic densities. The values rise almost exactly as $m^2$, as expected from the built-in $\langle\sigma v\rangle\propto m^{-2}$, with $x_f\approx24.4$--24.8 and $g_*=86.25$. Linear interpolation crosses 0.120 at $\approx1158$ GeV, with $\pm6$ GeV for the tool's $\pm0.0012$ band and about $\pm5$ GeV for $\pm0.001$. Direct calls give 0.120 at 1159 GeV and 0.1202 at 1160 GeV.

\begin{table}[h]
\centering
\begin{tabular}{cccc}
\toprule
$m_\chi$ (GeV) & $\langle\sigma v\rangle$ ($10^{-26}$ cm$^3$/s) & $\Omega h^2$ & tool tension ($\sigma$) \\
\midrule
900 & 3.73 & 0.0731 & 39.1 \\
1000 & 3.03 & 0.0898 & 25.1 \\
1080 & -- & 0.1045 & 12.9 \\
1100 & 2.50 & 0.1083 & 9.8 \\
1159 & 2.25 & 0.1200 & 0.0 \\
1160 & -- & 0.1202 & 0.1 \\
1200 & 2.10 & 0.1285 & 7.1 \\
1300 & 1.79 & 0.1503 & 25.2 \\
\bottomrule
\end{tabular}
\caption{Tool relic density with the higgsino flag set. The tension is against $0.12\pm0.0012$. Values at 1080 and 1160 GeV come from \texttt{check\_all\_constraints}.}
\label{tab:relic}
\end{table}

The tool's ``consistent with Planck'' flag reads True at every mass from 1000 to 1300 GeV despite tensions of 7--25$\sigma$, and False only at 900 GeV ($39\sigma$). This appears to be a bug.

With the higgsino flag off at 1100 GeV, $\Omega h^2 = 0.1342$, 0.1199 and 0.0909 for $\langle\sigma v\rangle = 2.0$, 2.25 and $3.0\times10^{-26}$ cm$^3$/s. Hence $\Omega\propto1/\langle\sigma v\rangle$ with $\Omega h^2\,\langle\sigma v\rangle\approx0.27\times10^{-26}$ cm$^3$/s, and the Planck abundance requires $2.25\times10^{-26}$ cm$^3$/s at 1.1 TeV.

Under the built-in scaling, a 10\% change in $\langle\sigma v\rangle$ moves the crossing mass by about 5\% (55--60 GeV). Matching the literature masses 1.08, 1.10 and 1.20 TeV would require effective factors of about 0.87, 0.90 and 1.07 on the tool's own rate. These are not physical Sommerfeld factors. The critic estimated that the Planck uncertainty shifts the mass by only $\sim\pm6$ GeV, negligible compared with the literature spread. The planned with/without-Sommerfeld comparison could not be done because the tool has no such switch.

\subsection{Inelastic LZ counts and the kinematic edge}
At $m_\chi=1100$ GeV the predicted LZ count falls steeply with splitting:
\begin{itemize}
\item $\delta=0$: $N_{\rm LZ}=1.10\times10^7$;
\item $\delta=200$ keV: $5.1\times10^4$;
\item $\delta=300$ keV: 744.6;
\item $\delta=330$ keV: 57.7;
\item $\delta=350$ keV: 1.97, with $P(N\ge1)=0.86$ and peak $E_R=315$ keV;
\item $\delta=355$ keV: 0.494;
\item $\delta=360$ keV: 0.085;
\item $\delta=365$ keV: 0.009;
\item $\delta=370$, 375, 380, 390, 500, 566 and 700 keV: 0. The spectrum peaks would sit at 333--630 keV.
\end{itemize}

The 248 keV event requires $v_{\min}$ of 762, 774, 781, 787 and 799 km/s at $\delta=350$, 360, 365, 370 and 380 keV. It is therefore accessible below the 776 km/s cutoff only up to $\delta\approx360$ keV. At $\delta=380$ keV, $v_{\min}(248\,\mathrm{keV})$ is 802, 799 and 796 km/s for $m=1000$, 1100 and 1200 GeV, so the edge depends only weakly on mass.

The hand estimate $\delta_{\max}=\mu v_{\max}^2/2$ gives $\approx368$ keV for 776 km/s and $\approx385$ keV for $544+250$ km/s, consistent with the tool's closure between 360 and 370 keV. The zero counts at $\delta\ge370$ keV are therefore produced by the hard velocity cutoff, not by the energy window. This explains why the first-pass constraint checks at $\delta=400$ keV returned $N_{\rm LZ}=0$, a result initially flagged as contradicting the literature.

Log-linear interpolation gives $\delta(N=1)=352.5$ keV and $\delta(N=0.3)=356.4$ keV at 1100 GeV, so the ``about one event'' band is only $\sim4$ keV wide. In the second pass the window was described as $\sim15$ keV wide in the broader sense of counts falling from 57.7 at 330 keV to zero by 370 keV. Along the thermal line the mass dependence is weak:
\begin{itemize}
\item at 1159--1160 GeV, $\delta=360$ keV gives $N=0.135$--0.136, compared with 0.085 at 1100 GeV;
\item at 1160 GeV, $\delta=350$ keV gives 2.52;
\item $\delta(N=1)\approx353$ keV at 1160 GeV.
\end{itemize}

At $\delta=350$ keV, the checks at 1080, 1100, 1160 and 1200 GeV give $N_{\rm LZ}=1.79$, 1.97, 2.52 and 2.90. The annual modulation at 1100 GeV, 350 keV and $E_R=248$ keV has a fractional amplitude of 0.985 peaking on June 2; on September 1 the rate is 22.5\% of the peak.

Comparison with the literature fits gives a mixed picture:
\begin{itemize}
\item The 350 keV value of arXiv:2609.01504 agrees directly with the tool.
\item The 377 keV value of arXiv:2609.02608 lies 24.5 keV above the tool's $\delta(N=1)$. About 17 keV of that offset would be explained if that paper used a 794 km/s cutoff, which is an assumption. The remaining 5--8 keV is unexplained.
\item The 339 keV sideband optimum of arXiv:2609.26698 does not depend on normalization, so it is not a like-for-like comparison. With the tool's normalization such splittings give tens to hundreds of events.
\end{itemize}

Halo sensitivity could not be computed because the tool has no halo inputs. Analytically, $v_{\rm esc}=500$ or 600 km/s ($v_{\max}\approx732$ or 832 km/s) would move the edge to $\approx327$ or $423$ keV. The first estimate held that 377 keV needs a $v_{\max}$ only 1--2\% higher ($v_{\rm esc}\gtrsim555$ km/s). The critic argued that the relevant quantity is $v_{\min}$ at 248 keV, which requires $v_{\max}\gtrsim795$ km/s, about 2.5\% higher, or $v_{\rm esc}\approx563$ km/s for an assumed $v_E\approx232$ km/s. On this basis, bridging 352 and 377 keV by halo uncertainty is plausible but not demonstrated. All halo choices considered stay below the 566 keV solar bound.

\subsection{Constraint checker and scans}
\texttt{check\_all\_constraints} rated every inelastic point ``VIABLE''. This includes $\delta=200$ keV at 1100 GeV, where $N_{\rm LZ}=5.1\times10^4$, so the overall pass flag ignores direct detection.

Only $\delta=0$ was excluded on all three counts:
\begin{itemize}
\item LZ: $N_{\rm LZ}=1.1\times10^7$;
\item Fermi-LAT: $\langle\sigma v\rangle_0=2.5\times10^{-24}$ cm$^3$/s, against a limit of $2.3\times10^{-25}$ cm$^3$/s;
\item CMB: $p_{\rm ann}=3.4\times10^{-28}$, against a limit of $3.2\times10^{-28}$.
\end{itemize}

For $\delta>0$, the tool's present-day annihilation rate is $\sim2\times10^{-42}$ to $5\times10^{-41}$ cm$^3$/s, labelled ``p-wave + mass-splitting suppressed''. The assistant judged this inappropriate for an s-wave, Sommerfeld-enhanced higgsino with $\langle\sigma v\rangle\sim10^{-26}$ cm$^3$/s. The first-pass critic marked this finding contested.

The $17\times17$ scan reported 144 of 289 points viable and 9 predicting 0.3--5 LZ events, with best fit $(988\ \mathrm{GeV}, 350\ \mathrm{keV})$. The tool's own relic curve gives $\Omega h^2\approx0.087$ at that best fit, so the point is under-abundant.

The $15\times15$ scan reported 210 of 225 points viable, 15 predicting 0.3--5 events (apparently one row near 350 keV), and best fit $(1000\ \mathrm{GeV},350\ \mathrm{keV})$ on the scan edge. Because of the flag problems, these viability counts are not physical statements.

\subsection{Elastic SI placement}
LZ 4.2 t$\cdot$yr (arXiv:2410.17036) states no 1 TeV value in its text. HEPData record 155182 gives, at 1008 GeV, an observed 90\% CL limit of $3.056\times10^{-47}$ cm$^2$ and a median sensitivity of $6.834\times10^{-47}$ cm$^2$, scaling roughly as $m^{0.985}$. Scaled to 1.1 TeV these become $\approx3.3\times10^{-47}$ (observed) and $\approx7.5\times10^{-47}$ cm$^2$ (median). The units were not checked against the figure.

XENONnT 3.1 t$\cdot$yr (arXiv:2502.18005) states $3.7\times10^{-46}\,\mathrm{cm^2}\times(M/1\,\mathrm{TeV})$ above 200 GeV, which gives $\approx4.1\times10^{-46}$ cm$^2$ at 1.1 TeV. In the first pass, from abstracts only, its minimum limit was $1.7\times10^{-47}$ cm$^2$ at 30 GeV, and $\sim10^{-45}$ cm$^2$ at 1 TeV was only an extrapolation. That earlier estimate was superseded.

For the first runs:
\begin{itemize}
\item XENONnT (arXiv:2303.14729) gives $6.08\times10^{-47}\,\mathrm{cm^2}\times(M/100\,\mathrm{GeV})$, about $6.7\times10^{-46}$ cm$^2$ at 1.1 TeV. That run used a more conservative median power constraint, and the improvement to the current result is about $1.6\times$.
\item LZ minimum limits went from $9.2\times10^{-48}$ cm$^2$ at 36 GeV (arXiv:2207.03764) to $2.2\times10^{-48}$ cm$^2$ at 40 GeV, about $4.2\times$. The first-run LZ value at 1 TeV was not obtained.
\end{itemize}

The pure-higgsino loop-level cross section is only bounded. Hill and Solon give $\sigma_{\rm SI}\lesssim10^{-48}$ cm$^2$ at 95\% CL, independent of $M$ in the heavy-WIMP limit (arXiv:1309.4092 Eq.~5; arXiv:1409.8290). No central value is given because of destructive spin-0/spin-2 interference; for comparison the pure wino is $1.3^{+1.2}_{-0.5}{}^{+0.4}_{-0.3}\times10^{-47}$ cm$^2$ in arXiv:1309.4092. Chen and Hill (arXiv:1912.07795) keep it at or below $10^{-48}$ cm$^2$ in the TeV regime. An NLO MSSM study (arXiv:2607.20588) finds cancellations that can push it below $10^{-50}$ cm$^2$.

Taking the xenon SI floor as $5.0\times10^{-48}$ cm$^2$ at 1 TeV, about $5.5\times10^{-48}$ cm$^2$ at 1.1 TeV (O'Hare, arXiv:2109.03116 data file), the higgsino at 1.1 TeV lies:
\begin{itemize}
\item at least $\sim33\times$ below the LZ observed limit;
\item at least $\sim75\times$ below the LZ median sensitivity;
\item at least $\sim400\times$ below XENONnT;
\item at least $\sim5.5\times$ below the floor, i.e.\ inside the neutrino fog.
\end{itemize}
All ratios are lower bounds, because the higgsino value is itself an upper bound. The LZ limit is itself only about $6\times$ above the floor at 1 TeV. XLZD (arXiv:2410.17137) was relayed as reaching $3\times10^{-49}$ cm$^2$ evidence potential at 40 GeV.

\subsection{Gaugino scale implied by the splitting}
Using $\delta_0\simeq m_Z^2(s_W^2/M_1+c_W^2/M_2)$ (arXiv:1410.4549) with $M_1=M_2=M$ and $m_Z^2=8315$ GeV$^2$ gives:

\begin{center}
\begin{tabular}{cc}
\toprule
$\delta$ (keV) & $M$ (GeV) \\
\midrule
350 & $2.38\times10^{7}$ \\
370 & $2.25\times10^{7}$ \\
385 & $2.16\times10^{7}$ \\
566 & $1.47\times10^{7}$ \\
\bottomrule
\end{tabular}
\end{center}

The xenon-accessible window thus needs gauginos at $\sim2.2$--$2.4\times10^4$ TeV, about three orders of magnitude above the 10--30 TeV assumed in the initial plan; the critic flagged that error. The solar bound requires $M\lesssim1.5\times10^4$ TeV. With wino exchange only, these values are about 22\% lower, and loop and running corrections (arXiv:1410.4549, arXiv:2609.09830) change them by tens of percent. A sub-agent reading of the same literature put the gaugino scale for 100--400 keV at roughly $2$--$8\times10^4$ TeV, with Higgs-exchange $\sigma_{\rm SI}\sim10^{-53}$ cm$^2$ there.

A tree-level Higgs-exchange estimate, $\sigma_{\rm SI}\sim5\times10^{-46}\,\mathrm{cm^2}\,(1\,\mathrm{TeV}/M)^2(1+\mathrm{sgn}\,\mu\sin2\beta)^2$ (from arXiv:2606.31013 and arXiv:1410.4549), reaches the LZ limit at $M\sim4$ TeV. That scale corresponds to splittings far above the inelastic window. The resulting claim that inelastic and Higgs-exchange elastic signals are mutually exclusive was marked contested by the critic.

\subsection{Indirect detection (abstracts only)}
The full text of arXiv:2405.13104 could not be downloaded, and all indirect-detection findings were marked contested. As relayed from abstracts:
\begin{itemize}
\item A 15-yr Fermi-LAT line-like search comes ``within an order one factor'' of detecting the higgsino for some Milky Way analog profiles, with H.E.S.S.\ subleading (arXiv:2405.13104).
\item The higgsino is ``out of reach by at least a factor of a few'' even at the ultimate sensitivity of current H.E.S.S.-like telescopes (arXiv:2210.03140).
\item The H.E.S.S.\ Inner Galaxy Survey (arXiv:2608.07234, 546 h) quotes a line limit of $2.3\times10^{-28}$ cm$^3$/s at 1 TeV for an Einasto profile and says the thermal higgsino is ``probed''. The abstract does not state exclusion.
\item No MAGIC-specific higgsino limit was found.
\end{itemize}
For future reach, CTA could be the first experiment sensitive to the thermal higgsino (arXiv:2008.00692). CTA and SWGO could detect the line at high significance for all but the most pessimistic profiles (arXiv:2405.13104). CTAO-North could reach the required sensitivity within a decade (arXiv:2506.08084).

\section{Discussion}
The combined picture is consistent across the tool and the literature. A thermal higgsino at $\sim1.1$ TeV is excluded as a Dirac-like state: $Z$-mediated elastic scattering overshoots LZ by about seven orders of magnitude. A neutral splitting is therefore mandatory. Below $\delta\approx330$ keV the inelastic rate is far too large.

Only a narrow band near 350 keV, bounded above by the halo velocity cutoff, produces of order one LZ event. That band lies below the IceCube solar-capture bound $\delta>566$ keV of arXiv:2609.02775. Splittings allowed by solar capture are beyond xenon's inelastic reach under any halo considered. In that regime only loop-level elastic scattering remains, at $\lesssim10^{-48}$ cm$^2$, which is inside the neutrino fog and at least $\sim33\times$ below current LZ sensitivity. The higgsino is therefore not excluded by direct detection, but current xenon experiments essentially cannot see it unless gaugino mixing raises the elastic rate.

The tool's quantitative content is limited. Its relic curve is a calibrated scaling with no Sommerfeld, coannihilation or bound-state physics, so the 1159 GeV crossing measures the calibration rather than higgsino physics. The \texttt{is\_higgsino} flag silently ignores a supplied $\langle\sigma v\rangle$. The Planck and overall-viability flags are unreliable, and the indirect-detection model is physically inappropriate for a higgsino. The halo is hard-coded, so the window's dominant uncertainty could not be tested. The tool also contains no published LZ limit curve to cross-check against. Its self-description attributes arXiv:2609.02608 to the LZ Collaboration, whereas metadata identifies that paper as a single-author interpretation by Di Mauro.

The literature side carries its own limitations. In the first pass, every arXiv-backed tool failed, so all citations rested on OpenAlex abstracts. In the second pass verbatim abstract citations succeeded, but full texts were never downloaded. The LZ 1 TeV limit came from HEPData rather than the paper text, and the \texttt{cite} quote check on arXiv:2410.17036 failed. Indirect-detection statements remain at abstract level, and the 1.08--1.2 TeV mass spread was not explained. The assistant's summary text also twice incorrectly claimed that no tool calculations had been performed.

\section{Conclusions}
\paragraph{Established in the session.}
\begin{itemize}
\item The cited literature places the thermal pure higgsino at about 1.08--1.2 TeV. Planck gives $\Omega_c h^2=0.120\pm0.001$.
\item The tool's calibrated relic curve gives $\Omega h^2=0.0731$--0.1503 over 0.9--1.3 TeV and crosses Planck at $\approx1159$ GeV.
\item In the tool, elastic $Z$ exchange is excluded. About one LZ event requires $\delta\approx352$ keV at a fixed 776 km/s cutoff, with closure by 370 keV. The literature solar-capture bound $\delta>566$ keV lies beyond that reach.
\item The loop-level elastic cross section, $\lesssim10^{-48}$ cm$^2$, lies at least $\sim33\times$ below LZ, $\sim400\times$ below XENONnT, and inside the neutrino fog at 1.1 TeV.
\item The xenon-accessible window corresponds to gauginos at $\sim2\times10^4$ TeV.
\end{itemize}

\paragraph{Contested or unresolved.}
\begin{itemize}
\item The mutual exclusivity of the inelastic and Higgs-exchange elastic signals.
\item The indirect-detection status, which is ``not excluded, profile-dependent, CTAO decisive'' at abstract level only
\end{itemize}

\section*{Data and code availability}
All outputs of this session are collected in the project folder \href{file:///Users/nramachandra/hep-genesis-testing/projects/ramp-combo-higgsino-dd}{\texttt{/\allowbreak{}Users/\allowbreak{}nramachandra/\allowbreak{}hep-genesis-testing/\allowbreak{}projects/\allowbreak{}ramp-combo-higgsino-dd}}.
Full documentation of how these results were obtained is in the supplementary material (\texttt{supplement.pdf}: implementation details, the tool-call log, the research board with every finding and its evidence, figures not shown here, and the code index) and the verbatim transcript (\texttt{transcript.pdf}), generated alongside this report.


\end{document}
